\chapter{1973:Ernst Otto Fischer and Geoffrey Wilkinson}

\label{chap:chemistry1973}

The Nobel Prize in Chemistry for the year 1973 was awarded jointly to Ernst Otto Fischer and Geoffrey Wilkinson.

\textbf{Motivation:}

for their pioneering work, performed independently, on the chemistry of the organometallic, so called sandwich compounds

\section{Ernst Otto Fischer}

\subsection*{Early Life and Education}

Ernst Otto Fischer was born on 1918-11-10 in Munich, Germany.

\subsection*{Career and Major Research}

Fischer studied at the Technical University of Munich, where he received his doctorate and later became a professor. He investigated the structure of ferrocene and other cyclopentadienyl complexes of the transition metals, and he discovered the transition-metal carbene complexes that bear his name.

\subsection*{Nobel Prize-winning Work}

The work for which Ernst Otto Fischer was awarded the Nobel Prize in Chemistry in 1973 was described by the Nobel Committee as: "for their pioneering work, performed independently, on the chemistry of the organometallic, so called sandwich compounds".

Working independently of Wilkinson, Fischer established by X-ray crystallography that ferrocene has a sandwich structure in which an iron atom is held between two cyclopentadienyl rings, and he extended this chemistry to other metals.

\subsection*{Significance and Impact}

Sandwich compounds opened a major field of organometallic chemistry with wide applications in catalysis and materials science.

\subsection*{Later Career and Honors}

Fischer remained at the Technical University of Munich for the rest of his career. He died on 2007-07-23 in Munich, Germany.

\subsection*{Legacy}

Metallocenes and related sandwich complexes are now fundamental building blocks in inorganic and organometallic chemistry.

\subsection*{Modern Developments}

Fischer's synthesis of coordination compounds and his understanding of metal-carbon bonds founded organometallic chemistry. The ferrocene sandwich compounds he and others explored became a cornerstone of organometallic catalysis, and modern homogeneous catalysis, including hydroformylation and cross-coupling reactions used in drug and materials synthesis, descends from his work.

\subsection*{Selected Publications}

"Zur Kristallstruktur des Ferrocens" (Zeitschrift für Naturforschung, 1952)

\clearpage

\section{Geoffrey Wilkinson}

\subsection*{Early Life and Education}

Geoffrey Wilkinson was born on 1921-07-14 in Todmorden, England.

\subsection*{Career and Major Research}

Wilkinson studied at Imperial College London, where he received his doctorate. After research positions in Canada and at Harvard University, he became professor at Imperial College. He established the sandwich structure of ferrocene and later introduced the rhodium complex known as Wilkinson's catalyst, widely used for hydrogenating alkenes.

\subsection*{Nobel Prize-winning Work}

The work for which Geoffrey Wilkinson was awarded the Nobel Prize in Chemistry in 1973 was described by the Nobel Committee as: "for their pioneering work, performed independently, on the chemistry of the organometallic, so called sandwich compounds".

Working with collaborators at Harvard, Wilkinson recognized that ferrocene, then called iron bis-cyclopentadienyl, consists of an iron atom sandwiched between two cyclopentadienyl rings.

\subsection*{Significance and Impact}

The discovery of the sandwich structure transformed the understanding of metal-carbon bonding and stimulated the growth of homogeneous catalysis.

\subsection*{Later Career and Honors}

Wilkinson remained at Imperial College until his retirement. He died on 1996-09-26 in London, England.

\subsection*{Legacy}

Wilkinson's catalyst remains a model for homogeneous hydrogenation, and metallocene chemistry is central to modern inorganic chemistry.

\subsection*{Modern Developments}

Wilkinson's work on organometallic chemistry, including the discovery and characterization of ferrocene, opened the modern era of homogeneous catalysis. The Wilkinson catalyst for hydrogenation and the countless transition-metal catalysts developed since are now indispensable in pharmaceutical synthesis, agrochemical production, and the manufacture of fine chemicals.

\subsection*{Selected Publications}

"The Structure of Iron Bis-Cyclopentadienyl" (Journal of the American Chemical Society, 1952)

\clearpage

\section*{Chapter Summary}

The 1973 prize recognized Fischer and Wilkinson for their independent work on sandwich compounds. Their elucidation of the structure of ferrocene founded the modern chemistry of metallocenes.

